%% Page formatting
\usepackage[T1]{fontenc}
%\usepackage[left=2cm,right=2cm,top=2cm,bottom=2cm]{geometry}

\usepackage{xspace}
\usepackage{xcolor}
\usepackage{tcolorbox}
\definecolor{light-gray}{gray}{0.85}

\usepackage{graphicx}

\usepackage{cleveref}
%\usepackage[hidelinks]{hyperref}

\usepackage{multirow}

%% Maths
%\usepackage{amsmath, amssymb, amsfonts}

%% Theorem styles
%% \usepackage{amsthm}
%% \theoremstyle{definition}
%% \newtheorem{definition}{Definition}
%% \newtheorem{conjecture}{Conjecture}
%% \newtheorem{theorem}{Theorem}
%% \newtheorem{lemma}{Lemma}

\usepackage{proof}

%% lstlisting and languages
\usepackage{listings}
\lstdefinelanguage{ContextualML}
{
  morekeywords={and, block, case, of, mlam, fn, impossible, let, in, schema,
    some, rec, type, ctype, prop, stratified, inductive, coinductive,
    LF, STLC,
    if, then, else,
    total, match, with,
    susp, force, down,
    ornament},
  keepspaces=true,
  sensitive,
  morecomment=[l]{\%},
  morecomment=[n]{\%\{}{\}\%},
  morestring=[b]"
}[keywords,comments,strings]

\lstloadlanguages{ContextualML}
\lstset{language=ContextualML}

% The order of the "literate" definitions is significant:
%   later definitions shadow earlier ones.  The \\Pi definition must come
%   *after* the \\ definition, or the first part of \\Pi --- that is, \\ --- will
%   be matched, and instead of $\Pi$ you'll get $\lambda Pi$.
%
\lstset{
  literate=
    {->}{{$\rightarrow~$}}2 %
    {=>}{{$\Rightarrow~$}}2 %
    {|-}{{$\vdash\,$}}2 %
    {..}{{$.\hspace{-0.025cm}.\hspace{-0.025cm}.$}}1 % is there any nicer way?
    {\\}{{$\lambda$}}1 %
    {\\Pi}{{$\Pi$}}1 %
    {\\gamma}{{$\gamma$}}1 %
    {\\Gamma}{{$\Gamma$}}1 %
    {\\Delta}{{$\Delta$}}1 %
    {\\psi}{{$\psi$}}1 %
    {\\sigma}{{$\sigma$}}1 %
    {\\Sigma}{{$\Sigma$}}1 %
    {\\Psi}{{$\Psi$}}1%
    {\\alpha}{{$\alpha$}}1%
    {\\beta}{{$\beta$}}1%
    {\\forall}{{$\forall$}}1%
    {\\.}{{$\cdot\,$}}1 %
    {FN}{{$\Lambda$}}1 %
    {[}{{$\lceil$}}1 % open box bracket
    {]}{{$\rceil$}}1 % close box bracket
    {u<}{{$\lfloor$}}1 % open unbox bracket
    {>u}{{$\rfloor$}}1 % close unbox bracket
    {<<b}{\color{blue}}1%
    {<<r}{\color{red}}1%
    {<<s}{\color{blue}}1%
    {<<d}{\color{red}}1%
    {<<m}{\color{green!60!black}}1%
    {>>}{\color{black}}1 %
    {?}{\bf{?}}1 %
    {UP}{{$\uparrow$}}1 % up shift
    {DOWN}{{$\downarrow$}}1, % down shift
  columns=[l]fullflexible,
        basicstyle=\ttfamily\lst@ifdisplaystyle\footnotesize\fi,
        keywordstyle=\bf,
        identifierstyle=\relax,
        stringstyle=\relax,
        commentstyle=\slshape\color{DimGrey},
        breaklines=true,
        % breakatwhitespace=true,   % doesn't do anything (?!)
        mathescape=true,   % interprets $...$ in listing as math mode
        xleftmargin=0.5cm,
      }



%% Acronyms
\usepackage[nolist,nohyperlinks]{acronym}
\newacro{PTS}{pure type system}
\newacro{MLTT}{Martin-L{\"o}f type theory}
\newacro{CoC}{calculus of constructions}
\newacro{OL}{object language}
\newacro{ML}{meta-language}
\newacro{NbE}{normalization by evaluation}
\newacro{HOAS}{higher-order abstract syntax}
\newacro{APTS}{adjoint pure type system}
\newacro{DSL}{domain-specific language}
\newacro{STLC}{simply-typed $\lambda$-calculus}
\newacro{DTLC}{dependently-typed $\lambda$-calculus}
\newacro{PER}{partial equivalence relation}

%% Bibliography
\usepackage{natbib}
\usepackage{bibentry}
\bibliographystyle{apalike}
\nobibliography*
